\documentclass[9pt]{extarticle}

\usepackage[utf8]{inputenc}
\usepackage[a4paper, margin=1.2cm]{geometry}
\usepackage[colorlinks=true, urlcolor=blue]{hyperref}
\usepackage{enumitem}
\usepackage{changepage}

\pagestyle{empty}

\setlist[itemize]{label=\textbullet, noitemsep, topsep=0pt}
\setlength{\parindent}{0pt}

% actual document shit

\begin{document}

\begin{center}
	{\huge\bfseries Paul Parisot}
	\\[4pt]
	{\Large Ingénieur Logiciel - Enseignement de l'Informatique}
	\\
\end{center}

% section for contact

\par\vspace{-5pt}
\rule{\linewidth}{0.4pt}
\par\vspace{-5pt}

\begin{center}
	Github: \href{https://github.com/paulogarithm}{paulogarithm}
	\\
	\href{mailto:paul.parisot@epitech.eu}{paul.parisot@epitech.eu}
	\\
\end{center}

\par\vspace{-10pt}
\rule{\linewidth}{0.4pt}
\par\vspace{7pt}

% section for the about us

{\large\bfseries À propos}
\par\vspace{-8pt}
\rule{\linewidth}{0.4pt}
\par\vspace{5pt}

Ingénieur logiciel et assistant technique (ASTEK) à Epitech pendant deux ans : {\bfseries plus de 400 étudiants} accompagnés en C, C++, assembleur et Haskell, avec un goût marqué pour la transmission.
\\
En préparation de l'agrégation d'informatique {\bfseries (algorithmique des graphes, complexité, automates)}, avec une formation complémentaire en management à McGill University, dont la sociologie du travail (Corporate Behaviour).
\\
Contributeur open-source (plus de 10 000 lignes de code sur des packages critiques) et pratique professionnelle de l'anglais (C1).
\par\vspace{5pt}

% section for skills

{\large\bfseries Compétences}
\par\vspace{-8pt}
\rule{\linewidth}{0.4pt}
\par\vspace{5pt}

\hspace*{5pt}{\bfseries Enseignés (ASTEK)}: C, C++, Assembleur, Haskell (dont multithreading)
\\
\hspace*{5pt}{\bfseries Maîtrise}: C/C++, Golang, Python, Lua(u), Haskell, OCaml, Assembleur, Shell, SQL
\\
\hspace*{5pt}{\bfseries Outils}: gdb, Valgrind, Ghidra, Git, Docker, Kubernetes, CI/CD, Environnement Linux
\\
\hspace*{5pt}{\bfseries Transverses}: pédagogie, création de contenus interactifs, animation de cours, travail en équipe
\\
\hspace*{5pt}{\bfseries Langues}: Français, Anglais (C1, professionnel)

\par\vspace{5pt}

% section for the professional experience

{\large\bfseries Expérience}
\par\vspace{-8pt}
\rule{\linewidth}{0.4pt}
\par\vspace{5pt}

{\bfseries Epitech / IONIS - Assistant Technique \& Formateur (ASTEK, 2 ans)}
\begin{itemize}[noitemsep, topsep=0pt]
	\item Enseignement à {\bfseries plus de 400 étudiants} en C, C++, asm et Haskell (dont des problèmes de multithreading) ; correction et évaluation des rendus.
	\item Création de {\bfseries contenus pédagogiques interactifs} (quiz, cours inversés, tutoriels) et contribution à un {\bfseries taux de réussite historiquement élevé} au sein des promotions coachées par l'ASTEK.
\end{itemize}

\par\vspace{5pt}

{\bfseries Gno.land - Ingénieur Logiciel Junior (Open-Source)}
\begin{itemize}[noitemsep, topsep=0pt]
	\item Développement de 6+ packages core en Gnolang (variante de Go pour smart contracts) ; correction de {\bfseries 6 bugs critiques} ; {\bfseries ~10k lignes de code} contribuées.
	\item Rédaction de tests exhaustifs sous des politiques strictes de CI/CD et de revue de code.
\end{itemize}

\par\vspace{5pt}

{\bfseries Faurecia / Forvia - Ingénieur Démo Logicielle}
\begin{itemize}[noitemsep, topsep=0pt]
	\item Conception d'un {\bfseries moteur de calcul d'itinéraire GPS} et d'une interface de recherche de POI, intégrés à une application démo présentée au {\bfseries CES 2024}.
	\item Livraison de fonctionnalités prêtes pour la production {\bfseries en avance sur le planning}, salué pour la rapidité et la qualité du code.
\end{itemize}

\par\vspace{5pt}

% personal projects

{\large\bfseries Projets Personnels}
\par\vspace{-8pt}
\rule{\linewidth}{0.4pt}
\par\vspace{5pt}

\underline{\bfseries Zappy -- Serveur TCP multi-clients en C}
\par
\begin{adjustwidth}{5pt}{0pt}
    Serveur TCP mono-thread à boucle événementielle ({\bfseries \texttt{select}}), codé seul : buffer borné par client et mise en quarantaine des clients qui le saturent ; aucune erreur ni fuite sous Valgrind ; tests unitaires et fonctionnels.
\end{adjustwidth}

\par\vspace{8pt}

\underline{\bfseries \href{https://github.com/paulogarithm/cclass}{cclass} -- Moteur de mutation mémoire à l'exécution}
\par
\begin{adjustwidth}{5pt}{0pt}
    Framework en C pur réécrivant à l'exécution des instructions assembleur en mémoire ({\bfseries \texttt{mmap}}, flags de protection, alignement strict) pour contourner le surcoût du dispatch par vtable.
\end{adjustwidth}

\par\vspace{8pt}

\underline{\bfseries \href{https://github.com/paulogarithm/zzzcript}{ZZZcript} -- Langage fonctionnel statiquement typé}
\par
\begin{adjustwidth}{5pt}{0pt}
    Lexer, parser et interpréteur AST from scratch, avec un moteur d'évaluation récursif en {\bfseries Go} inspiré de Haskell et Golang.
\end{adjustwidth}

\par\vspace{8pt}

% education

{\large\bfseries Formation}
\par\vspace{-8pt}
\rule{\linewidth}{0.4pt}
\par\vspace{5pt}

{\bfseries Epitech, Paris} \hfill 2022 -- 2027 \\
Diplôme d'ingénieur en technologies de l'information (diplôme de niveau Master) \\
Certification Niveau 7 - Expert en ingénierie logicielle

\vspace{6pt}

{\bfseries McGill University, Montréal} \hfill 2025 -- 2026 \\
Certificat en Management -- GPA : 3.57/4.0 \\
Cours suivis : Comptabilité, Managerial Economics \& Analysis, Corporate Behaviour (sociologie du travail), Business, Mathématiques, Gestion de projet, Marketing

\end{document}
